\documentclass[11pt,a4paper]{article}
\usepackage[T1]{fontenc}
\usepackage[utf8]{inputenc}
\usepackage{lmodern,amsmath,amssymb,mathtools}
\usepackage[margin=23mm]{geometry}
\usepackage{booktabs,longtable,array,enumitem}
\usepackage[hidelinks]{hyperref}
\usepackage{microtype}
\setlength{\parindent}{0pt}
\setlength{\parskip}{4pt plus 1pt}
\setlist{leftmargin=1.5em,itemsep=2pt,topsep=3pt}
\newcommand{\Set}{\mathbf{Set}}
\newcommand{\Cat}{\mathbf{Cat}}
\newcommand{\W}{\mathcal W_R}
\newcommand{\Aff}{\mathsf{AffCont}_{\mathrm{fin}}}
\newcommand{\id}{\mathrm{id}}
\newcommand{\one}{\mathbf 1}
\newcommand{\psint}{\int_{\mathrm{ps}}}
\newcommand{\status}[1]{\textbf{#1.}}
\title{O1 Feasibility and Scope Decision\\[5pt]
\large Independent compatibility with retained context and content grades}
\author{}
\date{1 November 2026}
\begin{document}
\maketitle
\vspace{-1em}

\status{Decision in brief}
The recommended O1 proof target is a genuinely indexed family: normalized
weighted contexts crossed with receipt grades over a finitary polynomial
cartesian monad. Its inputs are specified independently; pullback pasting
derives the key content-only reassignment comparison, with an exact bounded
recognition result for that recipe. Both grade categories are active and
nonthin, and continuation contexts genuinely depend on preceding results.
This supports a proof attempt, not acceptance of the complete
category-valued law or the full O1 gate.

The constant affine-container construction provides a demanding control:
its local cell and dependence tests pass, and duplication gives a sharp
bounded obstruction. It led to the stronger indexed route but is not being
substituted for that route. The ordinary affine/commutative-semiring mechanism
is already established mathematics. The remaining work is an explicit
retained-category lift, its full weak coherence, and a checked comparison
with that baseline, not discovery of independent traversal.

\status{Authority and stopping point}
This note answers the separately authorized focused O0/G1 feasibility
investigation under the provisionally adopted expanded contract~\cite{Contract}.
The source contract still bears its original ``proposed'' status; later
provisional adoption does not authorize its entire execution. The present
output is a theorem target, assumption ledger, primary-source comparison,
and worked tests. It is not G1 approval, a complete O0 charter for O1--O8,
O1 acceptance, or permission to proceed to applications. O1 remains decisive:
failure to obtain the required independent criterion cannot be repaired by
producing the later manuscripts faster. Fuller applications remain later
in the programme.

\section{What is inherited and what is still at issue}

The inherited project manuscript~\cite{Unified} supplies the relative ambient,
all-cell endpoint realization, and the weak pseudocoend/comma equivalence.
Its mixed factory constructs the algebra \emph{given} the complete coherent
crossing. Its finite probability/list law is a genuine positive control, not
the broader independent criterion. None of these results is relabelled new.

The strongest direct external baseline found is Fujii--Tsai--Montacute--Hasuo,
Theorem~42~\cite{FTMH}: a $\mathbf W$-operad and a
$\mathbf W$-commutative effect monad induce a Set-level law; compatibility
with the effect multiplication additionally requires ordinary commutativity
or a unary restriction. For finite injections, Proposition~36 identifies the
condition with commutativity plus affinity, so commutativity is available.
Proposition~35 handles heterogeneous coordinates; Examples~5 and~38 cover
normalized semiring effects. Their construction uses an ordinary coend,
not a category of retained sample witnesses.

This is a real independent mechanism, not a supplied-law factory. Matching
it to the present monoidal category of container grades still needs an
explicit bridge: their fixed operadic monad is not automatically our entire
grade category. Inferring output-level compatibility with shape maps from
ordinary naturality does not supply the native grade-cell, endpoint,
pseudobalancing, or weak-coherence lift. The audit establishes neither an
impossibility nor a priority claim for that residual result.

Earlier and complementary controls are also relevant. Finitary-container
traversal and commutative Kleisli compatibility are established in
\cite[Theorem~3.7 and Lemma~4.1]{JR}; normalized semiring distributions in
\cite[Example~3.3]{Fritz}. Traversable-functor representation
\cite[Theorems~5.6 and~5.8]{JO} concerns polymorphism over applicatives, not
arbitrary retained realizations. Polynomial distributivity from a terminal
universal problem~\cite[Proposition~2.2.3]{Weber} is independent but does not
itself produce a law between two pre-existing multiplications. Likewise
factorization recognition~\cite[Theorem~3.8]{RW} cannot solve the gluing
problem if its already assembled ambient is assumed. Monadic-container law
presentations~\cite{PD} are useful controls, not an automatic solution for
arbitrary independently specified grades.

\section{The fully specified constant-index candidate}

Fix nested universes large enough for the following categories and relative
pseudocoends; take $C=\Set$ in the smallest chosen universe. Shapes and
interface sets are small there; grade categories and Cat-valued tensors are
formed in larger universes. All pure set maps are allowed.

\subsection{Independent context data}
Let $R$ be a commutative semiring with $0\ne1$. An object of $\W$ is a finite
set $X$ with $p:X\to R$ and $\sum_xp(x)=1$. An arrow
$a:(X,p)\to(Y,q)$ is a function satisfying
\[
 q(y)=\sum_{a(x)=y}p(x).
\]
Zero-weight points are retained. This avoids an unwarranted closure claim
for full-support carriers over semirings with zero divisors. Put
\[
 K_A=\W^A,\qquad e_AP=\coprod_{a\in A}U(P_a),
\]
with substitution reindexing. Units are singleton samples. Dependent
accumulation has carrier $\coprod_{x\in X}Y_x$ and weight $p(x)q_x(y)$.
Independent tensor has product carrier and product weight, including the
singleton empty product. Normalization and closure of pushforward maps
follow from finite semiring distributivity. These monoidal sample tensors
are not asserted to be categorical products in $\W$.

\subsection{Independent content data}
An object $Q=(S,E)$ of $\Aff$ is a set $S$ of shapes with a finite position
set $E_s$ at each shape. Write
\[
 F_Q(B)=\coprod_{s\in S}B^{E_s}.
\]
An arrow $\beta=(h,k):Q\to Q'$ has a shape map $h:S\to S'$ and injective
backward maps $k_s:E'_{h(s)}\hookrightarrow E_s$. Its extension sends
$(s,b)$ to $(h(s),b k_s)$. Shape merging and occurrence deletion are allowed;
the two binary-to-unary projections are distinct parallel native cells.

Ordinary container substitution $Q\circ T$ has shapes $(s,t:E_s\to S_T)$
and positions $\coprod_{e\in E_s}E^T_{t(e)}$, so
$F_{Q\circ T}\cong F_QF_T$. Substitution of affine arrows is affine:
the backward leaf-address map is a composite of injections. Its unit $I$
has one shape and one position. Set
\[
 L_B=\Aff,\qquad d_BQ=F_Q(B),\qquad v_!Q=Q.
\]
In the inherited execution-order convention, $Q\diamond T=T\circ Q$,
so the content accumulation map has source $F_T(F_QB)$.
This convention fixes the order; no two opposite meanings of composition
are used. These data are specified without a mixed law.

\status{Exact scope}
The objects and morphisms of $K_A$ and $L_B$ vary independently, but
$L_B$ is constant in $B$. This is not a theorem for arbitrary indexed
content fibres, arbitrary contextads, measurable/QBS dependence, or
higher-order effects.

\section{The control target and common proof obligations}

\status{TARGET A constant indexed control not an established theorem}
For every such $R$, form the inherited modules
\[
 M_R(A,B)=(e_A\downarrow B),\qquad
 N(A,B)=(A\downarrow d_B),\qquad
 (X\odot_CY)(A,D)=\psint^{B\in C}X(A,B)\times Y(B,D).
\]
The specified data should derive, with no interaction supplied as input,
a complete pseudodistributive law
\[
 \tau_R:N\odot_CM_R\Longrightarrow M_R\odot_CN.
\]
The intended components, comparison directions, and acceptance conditions
are the following.

\begin{enumerate}
\item For $Q=(S,E)$ and $P\in K_B$, the generic crossing is
\begin{align*}
 \widehat P_{(s,b)}&=\bigotimes_{e\in E_s}P_{b(e)}\in\W,
 &\widehat Q&=Q\in L_{e_BP},\\
 \chi_{Q,P}:e_{F_QB}\widehat P&\longrightarrow F_Q(e_BP),
 &(s,b,(x_e))&\longmapsto(s,e\mapsto(b(e),x_e)).
\end{align*}
On a content cell $(h,k)$ and a context cell $a_b:P_b\to P'_b$, its
witness map is $(x_e)\mapsto(a_{b(k(e'))}(x_{k(e')}))_{e'}$.
\item These maps must define a functor on every native cell, actual
pseudobalancing data and their relations, and a module transformation in
both endpoints. They must retain the changing-endpoint squares
\[
 \alpha:P\to u^*P',\quad \beta:v_!Q\to Q',\qquad
 d_{B'}\beta\,d_v(Q)f=f'e_u(P')e_A\alpha,
\]
for all $u:A\to A'$, $v:B\to B'$, with both pasting operations.
\item Derive invertible comparisons with the four boundaries
\begin{align*}
 \tau(j_NM)&\cong Mj_N,&\tau(Nj_M)&\cong j_MN,\\
 \tau(\mu_NM)&\cong(M\mu_N)(\tau N)(N\tau),&
 \tau(N\mu_M)&\cong(\mu_MN)(M\tau)(\tau M).
\end{align*}
Prove the \emph{complete} coherence in the inherited definition of law,
including its ambient associators and interchangers: the Marmolejo
nine-equation presentation or the stated equivalent modern presentation,
not just these four boundaries. Then transport through the inherited
pseudocoend/comma equivalence to the relative algebra on
$H_R(A,B)=(e_A\downarrow d_B)$.
\item The theorem should be natural under unital semiring homomorphisms
$\varphi:R\to R'$, acting by $(X,p)\mapsto(X,\varphi p)$ and retaining
functions. Derive a pure-compatible strong comparison
$H_R\to H_{R'}$ with coherent identity and composite comparisons.
Finite sums and products commute with $\varphi$; no fullness or equivalence
is promised. The nonidentity map $\mathbb Q_{\ge0}\to\mathbb B$ taking
positive weights to true is the required substantive coefficient change,
from quantitative histories to possibilities, not a relabeling.
\item Opposite-side trivialization must recover each one-sided algebra.
Within this fixed mixed carrier the unit inclusions should also be locally
fully faithful: $I$ has only its identity endomorphism and singleton context
families have unique self maps, so the mixed hom equations reduce to the
one-sided equations. Prove their compatibility with composition separately.
Compare the inherited terminal-grade List control by a declared
monoid-object specialization using singleton and flattening maps; simply
fixing $Q=\mathrm{List}$ is not closed under substitution and is not that proof.
\item Within the prescribed independent-occurrence crossing and deterministic
weight-preserving witness maps, classify exactly the allowable ordinary
container maps. Under a normalized pair $c+d=1$, $cd\ne0$, they are
precisely the affine maps. The scope and proof of this bounded part are
below; it is not recognition of arbitrary unmarked realizations.
\end{enumerate}

\status{What is already calculated}
The displayed crossing is typed. Its local cell map preserves weights by
summing omitted normalized factors and applying each pushforward equation.
For a pure $v:B\to C$, both representatives have fibre
$\bigotimes_eP_{v(b(e))}$ and the same occurrence-addressed body.
Even noninjective $v$ changes labels without merging samples. Thus the
pure-middle generator has a concrete identity/reindexing comparison.

The proposed context-multiplication comparison is the addressed bijection
\[
 \prod_e\left(\coprod_{x_e\in P_e}T_{e,x_e}\right)
 \;\cong\;
 \coprod_{(x_e)\in\prod_eP_e}\prod_eT_{e,x_e},
\]
with common weight $\prod_e p_e(x_e)t_{e,x_e}(y_e)$. The content comparison
regroups a product over leaf addresses $(e,f)$ into nested products. Units
insert or remove singleton address levels. These calculations explain the
intended proof mechanism. A global definition and explicit verification of
every stipulated law diagram, its native-cell naturality, and pseudocoend
descent remain proof obligations; they have not been certified by this note.

\section{Two proved discriminating tests}

\subsection{A dependent positive witness with both grades active}
Use rational probabilities. A prior fair bit $x$ selects a binary content
shape labeled $(b_0,b_1)$ when $x=0$, and a unary shape labeled $b_1$ when
$x=1$. Let $P_{b_0}$ be a fair bit and $P_{b_1}$ uniform on four points.
The accumulated context has eight points of mass $1/16$ in the first branch
and four of mass $1/8$ in the second. The continuation is a dependent sum
of products of different cardinalities, not one fixed tensor.

Map both content shapes to one unary shape, keeping the first binary
position or the sole unary position. On context grades, use identity at
$b_0$ and parity at $b_1$. The combined witness function is
\[
 (0,(r_0,r_1))\longmapsto(0,r_0),\qquad
 (1,(r_1))\longmapsto(1,r_1\bmod2).
\]
For a target $(0,r)$, four source points contribute $4/16=1/4$.
For $(1,r)$, two contribute $2/8=1/4$. This proves the pushforward equation
for all four target points. Its retained-output equation follows by the
chosen projection and parity. It therefore tests real result dependence,
shape merging, occurrence deletion, and a noninvertible context map.

As an endpoint control, identify $b_0$ and $b_1$ by a pure map while keeping
a binary shape. Two fair occurrences still have four histories, not two.
First and second projection cells differ on $(0,1)$ even after label
identification. Retained cell distinctions are not erased by pure balance.
An exact rational checker independently confirms this witness and the
finite dependent-reassociation and Boolean controls; that computation is
a regression test, not proof of TARGET A.

\subsection{An exact obstruction to duplicated positions}
Allow an ordinary container arrow $(h,k)$ whose backward position maps need
not be injective, but keep the displayed crossing and its retained-evaluation
equation. Suppose $k(e'_1)=k(e'_2)=e$ with $e'_1\ne e'_2$. Choose $B=\one$
and the normalized two-point context with weights $c,d$, $cd\ne0$, at every
source occurrence. The body equation forces both target sample values to
equal $x_e$. Hence the pushforward assigns weight zero to the event
$(x_{e'_1},x_{e'_2})=(0,1)$. The target independent product assigns weight
$cd\ne0$, after marginalizing its other coordinates. No deterministic
weight-preserving lifting with that body equation exists.

Conversely, an injective $k$ retains distinct coordinates; summing each
omitted coordinate contributes $1$, which proves the required pushforward
equation for every normalized context family. Thus injectivity is necessary
and sufficient in this bounded class. For a fair bit the diagonal equality
event has probability $1$ instead of $1/2$; over Boolean full-support
weights the diagonal omits off-diagonal possibilities.

The nondegeneracy condition is essential for the universal necessity claim.
Natural-number normalized weights, for example, are only Dirac laws.
Pure Set-family contexts admit diagonals because there is no independent
weight-preservation condition. Correlated witnesses or explicit sharing
can also evade the obstruction by changing the preserved interface or
semantics. The result rules out neither every law nor every repair.

\section{The genuinely indexed content boundary}
\label{sec:indexed}
\status{Recommended TARGET B not an established theorem}
Fix independently the semiring $R$ above and a finitary polynomial cartesian
monad $(T,\eta,\mu)$ on $\Set$. Thus $T(B)=\coprod_{s\in S}B^{E_s}$
with finite position sets, $T$ preserves pullbacks, and the naturality squares
of $\eta$ and $\mu$ are pullbacks. These are content-only hypotheses; they
mention neither weighted contexts nor a mixed law. Unit and multiplication
carry canonical position bijections. List and the identity monad are controls;
the family is parameterized by $T$, not a disjoint choice of unrelated laws.

Define genuinely indexed content grades and their extension by
\[
 L^T_B=\Set/(TB),\qquad d_B(q:Q\to TB)=B+Q,\qquad
 \nu_Q=[\eta_B,q]:B+Q\longrightarrow TB.
\]
Native cells are maps $a:Q\to Q'$ with $q'a=q$; they can merge receipt tags
and have distinct parallel noninvertible maps. Pure $f_!$ postcomposes with
$Tf$, whereas $f^*Q=Q\times_{TB}TX$. The unit grade is empty. For
$r:Z\to T(B+Q)$, accumulation has carrier $Q+Z$ and label
\[
 [q,\,\mu_BT(\nu_Q)r]:Q+Z\longrightarrow TB;
\]
its extension map is coproduct reassociation. The monad identities prove the
object-level unit and associativity label equations without any sampling.

\status{Proved local structural comparison}
For $f:X\to B$, let $Q^f=Q\times_{TB}TX$ and
$\bar f:X+Q^f\to B+Q$. Cartesianness of $\eta$ and disjoint coproducts give
\[
 X+Q^f\ \cong\ (B+Q)\times_{TB}TX
 \quad\text{over }\nu_Q\text{ and }Tf.
\]
Apply $T$ to this pullback and paste the cartesian naturality square of
$\mu$. Pulling back along $r$ identifies the second summands in
\[
 f^*(Q\diamond Z)\ \cong\ Q^f\diamond(\bar f)^*Z.
\]
The first summands are both $Q^f$. This proves the canonical comparison on
grades. Its maps are induced by the pullback universal properties, not
chosen by assuming a weighted interaction. In the List case its inverse
splits a lifted flat list into the blocks fixed by the original $r$-list;
repeated receipts use distinct blocks and nullary receipts use empty blocks.
A finite exact check covers 55 such witnesses, including repetition and
nullary tags. Neither this lemma nor those tests verify every pseudo-law.

\status{Exact crossing target}
For $P\in K_B$, write $\pi:E=e_BP\to B$. For $y\in B+Q$, express
$\nu_Q(y)=(s,b:E_s\to B)$ and put
\[
 \widehat P_y=\bigotimes_{e\in E_s}P_{b(e)},\qquad
 \widehat Q=Q\times_{TB}TE\longrightarrow TE.
\]
The displayed pullback identifies the crossing carrier with
\[
 e_{B+Q}\widehat P\ \cong\ (B+Q)\times_{TB}TE
 \ \cong\ E+\widehat Q=d_E\widehat Q.
\]
TARGET B is that these formulas derive a complete law
$\tau_{R,T}:N_T\odot M_R\Rightarrow M_R\odot N_T$ and the induced
comma algebra, with all the native-cell, endpoint, balancing and full
coherence requirements of TARGET A, not just a bijection of carriers.
Slice pullback/postcomposition Beck--Chevalley handles the independent
endpoint comparison; weighted dependent sum/product distributivity supplies
the context comparison. Their global coherence remains to be proved.

A direct indexed-cell control uses two receipt tags $a,c$, both labeled
$[b,b]$, merged into one tag $m$, and a context cell from a uniform four-point
sample to a fair bit by parity. Each decorated source tag has 16 equiprobable
sample pairs; coordinatewise parity gives four equiprobable target pairs.
Thus both grade maps are active and noninvertible. There is no probability
measure on the tag set itself. The full crossing square is checked on
36 source-extension points, including the four ordinary outputs.

The declared parameter morphisms fix $T$ and vary $R$ by unital semiring
homomorphisms; no $T$-changing functoriality is claimed. The comparison under
$R\to R'$ is again required. The one-sided
unit inclusions are to be proved strong and locally fully faithful, now
using the empty receipt grade. A precise List comparison is available as a
target: send $(P,Q,f)$ to $(P,\nu_Qf)$ in the fixed-$T$ history semantics.
Prove this is a pure-compatible strong comparison; for $T=\mathrm{List}$
it compares with the inherited control. It forgets receipt distinctions
and is not asserted to be an equivalence.

\status{Proved bounded recognition of this recipe}
Assume $T$ is polynomial and preserves pullbacks, but do not initially assume
its unit or multiplication cartesian. In the canonical slice-receipt recipe,
the occurrence-display map is an isomorphism for every $f,Q$ and the
\emph{canonical tag-preserving} substitution comparison is an isomorphism
for every $f,Q,Z$ if and only if $\eta$ and $\mu$ are cartesian.
Sufficiency is the preceding pullback proof. For necessity, $Q=\varnothing$
makes the occurrence-display square exactly the naturality square of $\eta$.
For multiplication take $Q=TB$, $q=\id$, $Z=TTB$ and $r=T(\mathrm{inr})$.
Then $Q^f=TX$ and, because $T$ preserves the coproduct-injection pullback,
$(\bar f)^*Z\cong TTX$. The comparison on the $Z$-summand is exactly
\[
 TTX\longrightarrow TTB\times_{TB}TX,\qquad
 w\longmapsto(TTf(w),\mu_X(w)).
\]
Its bijectivity says that the naturality square of $\mu$ is a pullback.
Arbitrary coproduct isomorphisms would not justify this argument; preserving
the two tags and using the canonical comparison are essential.

This recognizes the content-only reassignment mechanism, not all possible
mixed laws. The full proof must connect it explicitly to the construction
and its preserved interface. The affine iff is a separate cell boundary,
not a converse for unrestricted indexed realizations.

The structural assumption has a hard boundary. For the finite-multiset
monad and $f:\{0,1\}\to\one$, the two nested bags
$\{\{0\},\{0,1\}\}$ and $\{\{1\},\{0,0\}\}$ have the same flattened
bag and the same image under $TTf$, but are distinct. Thus the naturality
square of its multiplication is not a pullback. This excludes that
unique-reassignment criterion, not every distributive law for bags.

Finally, TARGET B's slice cells preserve the complete $T$-label. They do
\emph{not} automatically include affine deletion or permutation cells from
TARGET A. Enlarging them requires an occurrence category with explicit
backward position maps; an ordinary natural $T$-valued label equality cannot
silently supply that extension. Repeated labels denote independently sampled
occurrences, not shared or cached samples.

\section{Assumption ledger and recommendation}

The following ledger separates entry data from conclusions. ``Open'' means
work to be proved, not an additional compatibility hypothesis.
\begin{longtable}{@{}p{0.08\linewidth}p{0.22\linewidth}p{0.62\linewidth}@{}}
\toprule
ID & Role and status & Statement, use, and boundary \\
\midrule
\endhead
A1 & Shared typing, supplied & Set base, nested universes, finite positions and
finite witness fibres. Required for products, sums and relative tensors;
no measurable or higher-order closure follows. \\
A2 & Shared weights, supplied & Commutative semiring with $0\ne1$; normalized
weights, raw carriers and pushforward-preserving maps. Normalization permits
deleting samples; commutativity permits permutation. Zero points are kept. \\
A3 & A only, supplied & Affine containers, with injective backward position
maps and arbitrary shape maps. Defined before any interaction; the duplicated
position witness explains the boundary. \\
D1 & Shared and A, derived & Normalization of products and dependent
sums, closure of affine substitution, and their address comparisons follow
from the independent definitions. Not a supplied mixed law. \\
D2 & A local, derived & Typed crossing, projection cell maps,
body equations, finite pushforward calculations, and pure-middle generator
comparison. The two worked tests establish their stated finite claims. \\
A4 & A necessity only & Existence of $c+d=1$ with $cd\ne0$ is used
only to exclude duplication. Sufficiency does not need it. \\
O1a & Both laws, open & Global pseudobalanced functor, endpoint action,
all unit and multiplication modifications, and complete law coherence.
The target conclusion must not be inserted here as an entry condition. \\
O1b & Both controls, open & Coefficient-change comparisons; strong one-sided
inclusions; explicit List comparison; baseline-to-retained-category bridge.
No inferred equivalence after forgetting weights or witnesses. \\
A5 & B only, supplied & TARGET B requires a finitary polynomial cartesian monad $T$; its slice recipe derives the local pullback comparison. The multiset witness marks a boundary. \\ 
D3 & B local, proved & Canonical tag-preserving reassignment comparison; its bounded iff characterization of cartesian $\eta,\mu$. No alternative-law necessity claimed. \\
O1c & B scope, open & TARGET B is genuinely indexed, not arbitrary contentads.
The full law and the connection from bounded recipe recognition to its
preserved interface remain open; TARGET A's affine iff is not their proof. \\
\bottomrule
\end{longtable}

\status{Recommended next decision}
Approve a full O1 proof attempt for TARGET B, the indexed polynomial-receipt
family, rather than freeze constant affine content as the entire answer.
Keep TARGET A as a tested boundary/control and a possible later cell
enlargement. This is one recommended route, not a menu of loosely related
targets. It has independently checkable algebraic inputs, genuine dependence
on both interfaces and grade maps, and a common occurrence mechanism.

The next proof dossier must close the global law and parameter comparisons,
connect the proved bounded recipe recognition to the law, and establish the one-sided
and List controls. Stop at a verified theorem with those exact commitments,
or at a precise failure requiring a scope decision. Do not convert local
success, an unchanged open obligation, or a simpler replacement theorem into
O1 acceptance. If the required independent criterion fails, report that the
full contract cannot be completed as written.

Propose retaining the original display/codisplay pullback/pushout programme
as an open comparison and obstruction branch, not calling it solved or
silently superseded. Neither that disposition nor TARGET B is adopted by
this recommendation. The authorized feasibility stage ends with this note
and its checked local evidence; proceeding to the full proof requires the
next scope decision.

\status{Acceptance boundary}
Even a successful proof of TARGET B would not by itself discharge the
entire contract. Boolean possibilities provide a nonstochastic mathematical
control, not an interpreted capability/resource application or the O6
no-recovery requirement. No language, QBS comparison, observation regime,
resource semantics, or application family is frozen by this note. Historical
novelty remains unverified and is not manufactured by a more elaborate
ambient; the contract permits an established criterion with a nontrivial,
explicit specialization.

\begin{thebibliography}{12}
\bibitem{Contract}
\emph{Expanded Research Contract Version 0.1}.
Project TeX source, 1 November 2026; provisionally adopted after drafting.
\bibitem{Unified}
\emph{One relative construction for context content and effects}.
Project TeX manuscript, 1 November 2026.
Relevant labels: \texttt{thm:coend}, \texttt{def:law},
\texttt{thm:mixed}, \texttt{thm:problaw}.

\bibitem{FTMH}
S. Fujii, Y. C. Tsai, Y. Montacute and I. Hasuo.
\emph{Monads and Distributive Laws in Substructural Contexts}, extended
version, 2026. Theorem 42; Propositions 35--36; Examples 5 and 38.
\url{https://arxiv.org/html/2605.13533v1}.

\bibitem{JR}
M. Jaskelioff and O. Rypacek.
\emph{An Investigation of the Laws of Traversals}, 2012.
\url{https://arxiv.org/abs/1202.2919}.

\bibitem{Fritz}
T. Fritz et al.
\emph{Representable Markov categories and comparison of statistical
experiments in categorical probability}, 2023. Example 3.3.
\href{https://pure.strath.ac.uk/ws/portalfiles/portal/163006824/Fritz_etal_TCS_2023_Representable_Markov_categories_and_comparison_of_statistical_experiments.pdf}{Author manuscript}.

\bibitem{JO}
M. Jaskelioff and R. O'Connor.
\emph{A Representation Theorem for Second-Order Functionals}.
\url{https://www.fceia.unr.edu.ar/~mauro/pubs/Representation.pdf}.

\bibitem{Weber}
M. Weber. \emph{Polynomials in categories with pullbacks}.
Proposition 2.2.3 and Theorem 4.1.4.
\url{https://arxiv.org/abs/1106.1983}.

\bibitem{RW}
R. Rosebrugh and R. J. Wood.
\emph{Distributive Laws and Factorization}. Theorem 3.8.
\url{https://mta.ca/~rrosebru/articles/fac.pdf}.

\bibitem{PD}
C. Purdy and S. Damato.
\emph{Distributive Laws of Monadic Containers}, CALCO 2025.
\href{https://drops.dagstuhl.de/storage/00lipics/lipics-vol342-calco2025/html/LIPIcs.CALCO.2025.4/LIPIcs.CALCO.2025.4.html}{Primary publication}.
\end{thebibliography}
\end{document}
